% 左图为构造函数，右图为两个等大群体的总体概率；不含观测数据。
\begin{tikzpicture}[figure picture,foundation picture]
 \path[use as bounding box] (0,-8) rectangle (168,53);
 \draw[foundation axis] (10,0)--(72,0) node[right] {$p$};
 \draw[foundation axis] (10,0)--(10,46);
 \node[anchor=south west] at (5,48) {事件频率};
 \draw[FigureBlue,line width=1.1pt] (10,0)--(68,42);
 \draw[FigureRed,line width=1.1pt,dashed]
  plot[domain=0:1,samples=41] ({10+58*\x},{42*\x*\x});
 \node[anchor=west] at (16,37) {理想校准};
 \node[anchor=west] at (44,9) {过度自信};
 \node[below] at (10,0) {$0$}; \node[below] at (68,0) {$1$};
 \node[left] at (10,42) {$1$};
 \draw[draw=FigureGrid,line width=.6pt] (85,-5)--(85,49);
 \node at (116,43) {$\pi_1=0.2$};
 \node at (150,43) {$\pi_2=0.8$};
 \node[anchor=east] at (100,26) {常数};
 \node[foundation node,fill=FigureBlueFill,minimum width=26mm] at (116,26) {$p=0.5$};
 \node[foundation node,fill=FigureBlueFill,minimum width=26mm] at (150,26) {$p=0.5$};
 \node[anchor=east] at (100,8) {分组};
 \node[foundation node,fill=FigureRedFill,minimum width=26mm] at (116,8) {$p=0.2$};
 \node[foundation node,fill=FigureRedFill,minimum width=26mm] at (150,8) {$p=0.8$};
\end{tikzpicture}
